\chapter{Topologi Garis Real}\label{ch:b1:topology}

Limit terus-menerus merujuk pada kosakata geometri yang sama: titik
yang ``dekat dengan'' sebuah \omterm{def:b1:logic:sets}{himpunan}, \omterm{def:b1:logic:sets}{himpunan} ``tanpa kebocoran \omterm{def:b1:topology:closure}{perbatasan}'', \omterm{prop:b1:reals:intervals}{selang} yang
di dalamnya barisan tak dapat kabur. Bab ini memakukan kosakata itu ---
\omterm{def:b1:topology:open}{himpunan terbuka} dan \omterm{def:b1:topology:closed}{tertutup}, \omterm{def:b1:topology:closure}{interior} dan \omterm{def:b1:topology:closure}{penutup}, \omterm{def:b1:topology:dense}{kepadatan} --- pada garis
real, lalu membuktikan kekompakan ruas dalam bentuk barisannya.
Gagasan yang sama, pada ruang vektor bernorma, adalah bahan tahun kedua;
sedangkan pada $\R$ ia terjangkau dan langsung berguna bagi
\cref{ch:b1:continuity}.

\section{Himpunan terbuka, himpunan tertutup}

\begin{definition}[Persekitaran, himpunan terbuka]\label{def:b1:topology:open}
Sebuah \omterm{def:b1:logic:sets}{himpunan} $V \subseteq \R$ disebut \emph{persekitaran}\index{persekitaran}
$x \in \R$ bila ia memuat \omterm{prop:b1:reals:intervals}{selang} $\intoo{x - r}{x + r}$ untuk
suatu $r > 0$. Sebuah \omterm{def:b1:logic:sets}{himpunan} $U \subseteq \R$ disebut \emph{terbuka}\index{himpunan terbuka}
bila ia merupakan persekitaran bagi setiap titiknya:
\[
\forall x \in U,\ \exists r > 0, \quad \intoo{x - r}{x + r}
\subseteq U .
\]
\end{definition}

\begin{example}\label{ex:b1:topology:open}
\omterm{prop:b1:reals:intervals}{Selang} \omterm{def:b1:topology:open}{terbuka} bersifat \omterm{def:b1:topology:open}{terbuka}: karena untuk $x \in \intoo{a}{b}$, ambillah $r = \min(x -
a,\, b - x) > 0$. Setengah garis $\intoo{a}{+\infty}$ \omterm{def:b1:topology:open}{terbuka}; adapun $\R$ dan
$\emptyset$ \omterm{def:b1:topology:open}{terbuka} (yang terakhir secara hampa). Sedangkan $\intcc{0}{1}$
\emph{tidak} \omterm{def:b1:topology:open}{terbuka}: karena tak ada \omterm{prop:b1:reals:intervals}{selang} di sekitar $0$ yang tinggal di dalamnya.
\end{example}

\begin{proposition}[Kestabilan himpunan terbuka]\label{prop:b1:topology:openstable}
Sebarang gabungan \omterm{def:b1:topology:open}{himpunan terbuka} bersifat \omterm{def:b1:topology:open}{terbuka}; dan irisan \emph{hingga} \omterm{def:b1:topology:open}{himpunan
terbuka} bersifat \omterm{def:b1:topology:open}{terbuka}. Adapun irisan tak hingga boleh gagal:
$\bigcap_{n \geq 1} \intoo{-\frac 1n}{\frac 1n} = \{0\}$, yang tak \omterm{def:b1:topology:open}{terbuka}.
\end{proposition}

\begin{proof}
Gabungannya: jika $x \in \bigcup_i U_i$, maka $x \in U_{i_0}$ untuk suatu $i_0$,
dan \omterm{prop:b1:reals:intervals}{selang} yang dipasok $U_{i_0}$ duduk di dalam gabungannya. Adapun irisan
yang hingga: jika $x \in U_1 \cap \dots \cap U_k$, ambillah $r = \min(r_1,
\dots, r_k) > 0$ atas jari-jari yang dipasok masing-masing $U_j$. Untuk
contoh penyangkalnya: sebarang \omterm{prop:b1:reals:intervals}{selang} di sekitar $0$ memuat suatu $\frac{1}{n}$
(menurut Archimedes), sehingga meninggalkan irisannya.
\end{proof}

\begin{example}[Mengesahkan keterbukaan dengan jari-jari yang eksplisit]\label{ex:b1:topology:radii}
Apakah $U = \{x \in \R : x^2 > 2\}$ \omterm{def:b1:topology:open}{terbuka}? Ya, dan sertifikatnya
dapat dituliskan: $U = \intoo{-\infty}{-\sqrt2} \cup
\intoo{\sqrt2}{+\infty}$, yaitu gabungan dua setengah garis yang \omterm{def:b1:topology:open}{terbuka}, sehingga \omterm{def:b1:topology:open}{terbuka}
menurut \cref{prop:b1:topology:openstable}. Sebagai alternatifnya, berargumenlah
titik demi titik: untuk $x \in U$ dengan $x > \sqrt 2$, ambillah $r = x -
\sqrt2 > 0$: maka setiap $y \in \intoo{x - r}{x + r}$ memenuhi $y >
\sqrt 2$, sehingga $y^2 > 2$; simetris dengan itu di sebelah kirinya. Kedua
gayanya penting --- yang struktural (yaitu membangun dari \omterm{def:b1:topology:open}{himpunan terbuka}
yang dikenal lewat gabungan dan irisan hingga) lebih baik penskalaannya, sedangkan yang
bergaya $\varepsilon$ berjalan ketika tak ada struktur yang kasatmata; dan
\cref{ch:b1:continuity} akan menambahkan yang ketiga, yang paling ampuh:
bahwa $U$ merupakan \omterm{def:b1:logic:map}{prapeta} dari $\intoo{2}{+\infty}$ yang \omterm{def:b1:topology:open}{terbuka} di bawah
$x \mapsto x^2$ yang kontinu.
\end{example}

\begin{definition}[Himpunan tertutup]\label{def:b1:topology:closed}
Sebuah \omterm{def:b1:logic:sets}{himpunan} $F \subseteq \R$ disebut \emph{tertutup}\index{himpunan tertutup} bila
komplemennya $\R \setminus F$ \omterm{def:b1:topology:open}{terbuka}. Menurut de Morgan dan
\cref{prop:b1:topology:openstable}: sebarang irisan \omterm{def:b1:logic:sets}{himpunan} tertutup
bersifat tertutup, dan gabungan hingga \omterm{def:b1:logic:sets}{himpunan} tertutup bersifat tertutup.
\end{definition}

\begin{theorem}[Pencirian himpunan tertutup lewat barisan]\label{thm:b1:topology:seqclosed}
\omterm{def:b1:logic:sets}{Himpunan} $F$ \omterm{def:b1:topology:closed}{tertutup} jika dan hanya jika: untuk setiap barisan $(u_n)$ berisi titik
$F$ yang konvergen ke suatu $\ell \in \R$, limitnya $\ell$ termasuk
$F$. (Jadi ``\omterm{def:b1:topology:closed}{tertutup}'' $=$ ``stabil terhadap limit''.)
\end{theorem}

\begin{proof}
($\Rightarrow$) Misalkan $F$ \omterm{def:b1:topology:closed}{tertutup}, $u_n \in F$, $u_n \to \ell$, lalu
andaikan $\ell \notin F$. Komplemennya \omterm{def:b1:topology:open}{terbuka}: sehingga suatu
$\intoo{\ell - r}{\ell + r}$ menghindari $F$. Padahal kekonvergenannya menaruh $u_n$
pada \omterm{prop:b1:reals:intervals}{selang} itu untuk $n$ yang besar: yang bertentangan dengan $u_n \in F$.

($\Leftarrow$) Andaikan $F$ tak \omterm{def:b1:topology:closed}{tertutup}: maka komplemennya tak \omterm{def:b1:topology:open}{terbuka},
sehingga suatu $x \notin F$ tak mempunyai \omterm{prop:b1:reals:intervals}{selang} $\intoo{x - r}{x + r}$ di dalam
komplemennya; lalu dengan mengambil $r = \frac{1}{n+1}$, pilihlah $u_n \in F$ dengan
$\abs{u_n - x} < \frac{1}{n+1}$. Maka $u_n \in F$, $u_n \to x \notin
F$: jadi sifat barisannya gagal.
\end{proof}

\begin{example}[Uji barisannya, dari kedua arah]\label{ex:b1:topology:seqtest}
\emph{\omterm{def:b1:topology:closed}{Tertutup}:} $F = \Z \cup \bigl\{n + \frac1n : n \geq 2\bigr\}$.
Misalkan $u_k \in F$ dengan $u_k \to \ell$. Jendela $\intcc{\ell -
1}{\ell + 1}$ hanya memuat titik $F$ yang berhingga banyak (yaitu bilangan bulat
yang berhingga banyak, dan $n + \frac1n$ yang berhingga banyak), lalu di luar suatu peringkat
semua $u_k$ terletak di dalamnya: sehingga barisannya mengambil nilai yang berhingga
banyak, dan karena konvergen, ia akhirnya konstan (seperti pada
\cref{exo:b1:topology:3}): jadi $\ell \in F$. \omterm{def:b1:topology:closed}{Tertutup} --- meskipun
$F$ memuat pasangan titik yang berjarak $\frac1n$, sedekat-dekatnya.

\emph{Tak \omterm{def:b1:topology:closed}{tertutup}:} $G = \bigl\{\frac1m + \frac1n : m, n \in
\N^*\bigr\}$. Barisan $\frac1n + \frac1n \in G$ menuju
$0$, padahal $0 \notin G$ (karena ia jumlah dua suku yang positif): jadi
uji barisannya gagal, sehingga $G$ tak \omterm{def:b1:topology:closed}{tertutup}. Menariknya, setiap
$\frac1m$ \emph{memang} termasuk $\overline G \cap G$: memang
$\frac1m = \frac{1}{m+1} + \frac{1}{m(m+1)} \in G$. Inti gagasan
penutupnya: bahwa untuk membuktikan ketertutupan, kendalikanlah \emph{semua} barisan
yang konvergen sekaligus (biasanya lewat kehinggaan lokal atau lewat argumen
rumus \omterm{def:b1:topology:closed}{tertutup}); sedangkan untuk membantahnya, satu barisan pelarian yang terpilih baik
sudah cukup --- ketaksetangkupan itulah yang membuat arah negatifnya
menjadi yang mudah, dan contoh penyangkal bab ini semuanya berbentuk
satu baris seperti ini.
\end{example}

\begin{example}\label{ex:b1:topology:closed}
Ruas $\intcc{a}{b}$, setengah garis $\intco{a}{+\infty}$, \omterm{def:b1:counting:card}{himpunan
hingga}, dan $\Z$ (karena barisan bilangan bulat yang konvergen akhirnya konstan)
bersifat \omterm{def:b1:topology:closed}{tertutup}. Sedangkan $\intoc{0}{1}$ tak \omterm{def:b1:topology:open}{terbuka} (karena gagal di $1$) dan tak \omterm{def:b1:topology:closed}{tertutup}
(karena $\frac 1n \to 0 \notin$ \omterm{def:b1:logic:sets}{himpunannya}): jadi kebanyakan \omterm{def:b1:logic:sets}{himpunan} terbukan keduanya. Adapun $\R$ dan
$\emptyset$ sekaligus \omterm{def:b1:topology:open}{terbuka} dan \omterm{def:b1:topology:closed}{tertutup} --- dan hanya keduanyalah
\omterm{def:b1:logic:sets}{himpunan} bagian $\R$ yang demikian (\cref{exo:b1:topology:9}).
\end{example}

\begin{example}[Himpunan terbuka yang dirakit dari kepingan tak hingga banyak]\label{ex:b1:topology:rminusz}
Di sini $\R \setminus \Z = \bigcup_{n \in \Z} \intoo{n}{n+1}$: yaitu
gabungan tak hingga atas \omterm{prop:b1:reals:intervals}{selang} \omterm{def:b1:topology:open}{terbuka}, sehingga \omterm{def:b1:topology:open}{terbuka} menurut
\cref{prop:b1:topology:openstable} --- jadi $\Z$ \omterm{def:b1:topology:closed}{tertutup} tanpa
memerlukan argumen barisan apa pun. Perhatikan pembagian kerja pada
aturan kestabilannya: bahwa \emph{gabungan} \omterm{def:b1:topology:open}{himpunan terbuka} boleh sebarang
(karena setiap titiknya hanya memerlukan sertifikatnya sendiri, yang dipasok satu
\omterm{def:b1:logic:sets}{himpunan} yang memuatnya), sedangkan \emph{irisannya} harus tetap hingga
(sebab sertifikatnya harus diiriskan, dan jari-jari yang tak hingga banyak dapat
menyusut menjadi nol). \Cref{exo:b1:topology:10} akan menunjukkan bahwa contoh
ini merupakan bentuk umumnya: bahwa setiap \omterm{def:b1:logic:sets}{himpunan} bagian $\R$ yang \omterm{def:b1:topology:open}{terbuka} merupakan
gabungan lepas yang terbilang atas \omterm{prop:b1:reals:intervals}{selang} \omterm{def:b1:topology:open}{terbuka}.
\end{example}

\section{Interior, penutup, kepadatan}

\begin{definition}[Interior, penutup, perbatasan]\label{def:b1:topology:closure}
Misalkan $A \subseteq \R$.
\begin{itemize}
  \item Sebuah titik $x$ disebut \emph{interior} bagi $A$ bila $A$ merupakan
        \omterm{def:b1:topology:open}{persekitaran} $x$; dan \emph{interiornya}
        $\mathring{A}$\index{interior} adalah \omterm{def:b1:logic:sets}{himpunan} titik interiornya.
  \item Sebuah titik $x$ disebut \emph{lekat} pada $A$ bila setiap \omterm{def:b1:topology:open}{persekitaran}
        $x$ berpotongan dengan $A$; dan \emph{penutupnya}
        $\overline{A}$\index{penutup} adalah \omterm{def:b1:logic:sets}{himpunan} titik lekatnya.
  \item Adapun \emph{perbatasannya} adalah $\partial A = \overline A \setminus
        \mathring A$.\index{perbatasan}
\end{itemize}
Maka $\mathring A \subseteq A \subseteq \overline A$.
\end{definition}

\begin{proposition}[Sifat utamanya]\label{prop:b1:topology:closureprops}
\begin{enumerate}
  \item $\mathring A$ merupakan \omterm{def:b1:topology:open}{himpunan terbuka} terbesar yang termuat $A$; dan $A$
        \omterm{def:b1:topology:open}{terbuka} jika dan hanya jika $A = \mathring A$.
  \item $\overline A$ merupakan \omterm{def:b1:topology:closed}{himpunan tertutup} terkecil yang memuat $A$; dan $A$
        \omterm{def:b1:topology:closed}{tertutup} jika dan hanya jika $A = \overline A$.
  \item (Pencirian kelekatan lewat barisan) Di sini $x \in \overline
        A$ jika dan hanya jika $x$ merupakan limit sebuah barisan berisi titik
        $A$.
  \item Pengomplemenan menukar kedua gagasannya: $\R \setminus
        \overline A = \bigl(\R \setminus A\bigr)^{\!\circ}$.
\end{enumerate}
\end{proposition}

\begin{proof}
(4) Di sini $x \notin \overline A$ $\iff$ suatu \omterm{def:b1:topology:open}{persekitaran} $x$ menghindari $A$
$\iff$ suatu \omterm{prop:b1:reals:intervals}{selang} di sekitar $x$ terletak di $\R \setminus A$ $\iff$ $x$
bersifat \omterm{def:b1:topology:closure}{interior} bagi $\R \setminus A$.

(1) \omterm{def:b1:logic:sets}{Himpunan} $\mathring A$ \omterm{def:b1:topology:open}{terbuka}: karena jika $x \in \mathring A$, maka suatu
$\intoo{x-r}{x+r} \subseteq A$; lalu setiap titik $y$ pada \omterm{prop:b1:reals:intervals}{selang} itu mempunyai
\omterm{prop:b1:reals:intervals}{selang} yang lebih kecil di sekitarnya yang berada di dalamnya, sehingga di dalam $A$: jadi seluruh
\omterm{prop:b1:reals:intervals}{selangnya} termasuk $\mathring A$. Adapun sebarang \omterm{def:b1:topology:open}{himpunan terbuka} $U \subseteq A$ terdiri atas
titik \omterm{def:b1:topology:closure}{interior} $A$, sehingga $U \subseteq \mathring A$: jadi terbesar. Lalu
pencirian keterbukaannya menyusul.

(2) Secara terperinci. Menurut (4), $\R \setminus \overline A$ merupakan \omterm{def:b1:topology:closure}{interior}
$\R \setminus A$, yaitu \omterm{def:b1:topology:open}{himpunan terbuka} menurut (1): sehingga $\overline A$
\omterm{def:b1:topology:closed}{tertutup}, dan ia memuat $A$. Keminimalannya: misalkan $F \supseteq A$
\omterm{def:b1:topology:closed}{tertutup}. Maka $\R \setminus F$ \omterm{def:b1:topology:open}{terbuka} dan termuat $\R
\setminus A$, sehingga menurut kemaksimalan pada (1),
\[
\R \setminus F \subseteq \bigl(\R \setminus A\bigr)^{\!\circ}
= \R \setminus \overline A ,
\]
lalu mengambil komplemennya lagi: $\overline A \subseteq F$. Jadi
$\overline A$ merupakan pemuat \omterm{def:b1:topology:closed}{tertutup} yang terkecil. Pencirian sifatnya:
jika $A = \overline A$ maka $A$ \omterm{def:b1:topology:closed}{tertutup} (sebagaimana baru ditunjukkan); sedangkan jika $A$
\omterm{def:b1:topology:closed}{tertutup}, ia sendiri merupakan pemuat \omterm{def:b1:topology:closed}{tertutup} bagi $A$, sehingga keminimalannya
memaksa $\overline A \subseteq A$, jadi sama.

(3) Jika $u_n \in A$, $u_n \to x$: maka setiap \omterm{def:b1:topology:open}{persekitaran} $x$ memuat
suatu $u_n \in A$, sehingga $x \in \overline A$. Sebaliknya, jika $x \in
\overline A$: maka setiap \omterm{prop:b1:reals:intervals}{selang} $\intoo{x - \frac{1}{n+1}}{x +
\frac{1}{n+1}}$ berpotongan dengan $A$ pada suatu $u_n$, dan $u_n \to x$.
\end{proof}

\begin{example}\label{ex:b1:topology:closureexamples}
Di sini $\overline{\intoo{0}{1}} = \intcc{0}{1}$;
$\mathring{\intcc{0}{1}} = \intoo{0}{1}$; $\partial\intoo{0}{1} =
\{0, 1\}$. Untuk $A = \{\frac 1n : n \in \N^*\}$: $\overline A = A
\cup \{0\}$, $\mathring A = \emptyset$, $\partial A = A \cup \{0\}$.
Untuk $\Q$: menurut \omterm{def:b1:topology:dense}{kepadatannya} (\cref{thm:b1:reals:density}) setiap bilangan real
lekat pada $\Q$, sehingga $\overline{\Q} = \R$ sedangkan $\mathring{\Q} =
\emptyset$ (karena setiap \omterm{prop:b1:reals:intervals}{selang} memuat bilangan irasional): jadi \omterm{def:b1:topology:closure}{perbatasan}
$\Q$ adalah seluruh $\R$.
\end{example}

\begin{example}[Anatomi yang lengkap]\label{ex:b1:topology:anatomy}
Misalkan $A = \intoc{0}{1} \,\cup\, \bigl(\Q \cap \intoo{2}{3}\bigr)
\,\cup\, \{4\}$. Kita menghitung ketiga \omterm{def:b1:logic:sets}{himpunan} pada
\cref{def:b1:topology:closure}, kepingan demi kepingan.

\emph{\omterm{def:b1:topology:closure}{Interiornya}.} Sebuah titik $\intoo{0}{1}$ mempunyai seluruh \omterm{prop:b1:reals:intervals}{selang}
di dalam $A$: jadi \omterm{def:b1:topology:closure}{interior}. Adapun titik $1$: setiap \omterm{prop:b1:reals:intervals}{selang} di sekitarnya
bocor ke kanan $1$, yang di situ $A$ tak punya apa-apa sampai $2$: jadi bukan
\omterm{def:b1:topology:closure}{interior}. Tak ada titik $\Q \cap \intoo{2}{3}$ yang \omterm{def:b1:topology:closure}{interior} (karena setiap
\omterm{prop:b1:reals:intervals}{selang} memuat bilangan irasional, \cref{thm:b1:reals:density});
demikian pula $4$ yang terasing. Jadi $\mathring A = \intoo{0}{1}$.

\emph{\omterm{def:b1:topology:closure}{Penutupnya}.} Limit titik $A$ adalah: seluruh $\intcc{0}{1}$
(karena $0 = \lim \frac1n$ dengan $\frac 1n \in A$); seluruh
$\intcc{2}{3}$ (karena setiap bilangan real di sana merupakan limit bilangan rasional pada
\omterm{prop:b1:reals:intervals}{selangnya}, menurut \omterm{def:b1:topology:dense}{kepadatan} lagi); dan $4$. Tak ada yang lain: karena titik
di luar $\intcc{0}{1} \cup \intcc{2}{3} \cup \{4\}$ berjarak positif
ke \omterm{def:b1:topology:closed}{himpunan tertutup} itu. Jadi $\overline A = \intcc{0}{1} \cup
\intcc{2}{3} \cup \{4\}$.

\emph{\omterm{def:b1:topology:closure}{Perbatasannya}.} Di sini $\partial A = \overline A \setminus \mathring A
= \{0, 1\} \cup \intcc{2}{3} \cup \{4\}$.

Inti gagasan penutupnya: bahwa ketiga operasinya bekerja secara \emph{lokal} ---
karena setiap kepingan $A$ menyumbang menurut sifatnya sendiri (sebuah
\omterm{prop:b1:reals:intervals}{selang} yang padu mempertahankan bagian dalamnya, kepingan yang \omterm{def:b1:topology:dense}{padat} tetapi berpori berubah
seluruhnya menjadi \omterm{def:b1:topology:closure}{perbatasan}, sedangkan titik yang terasing murni \omterm{def:b1:topology:closure}{perbatasan}), dan
gambaran dua baris atas $A$ meramalkan setiap jawabannya sebelum satu bukti pun
dituliskan.
\end{example}

\begin{remark}[Jebakan yang lazim dalam penalaran himpunan titik]
(i) \emph{``Tak \omterm{def:b1:topology:open}{terbuka}'' tak berarti ``\omterm{def:b1:topology:closed}{tertutup}''}: karena kebanyakan
\omterm{def:b1:logic:sets}{himpunan} terbukan keduanya ($\intoc{0}{1}$), sedangkan dua \omterm{def:b1:logic:sets}{himpunan} justru keduanya
($\emptyset$, $\R$) --- jadi \omterm{def:b1:topology:open}{terbuka} dan \omterm{def:b1:topology:closed}{tertutup} bukanlah lawan melainkan
dual lewat pengomplemenan. (ii) \emph{\omterm{def:b1:topology:closure}{Interior} dan \omterm{def:b1:topology:closure}{penutup}
tak berkomutasi}: karena untuk $A = \Q$,
\[
\overline{\mathring A} = \overline\emptyset = \emptyset
\qquad\text{sedangkan}\qquad
\bigl(\,\overline A\,\bigr)^{\!\circ} = \mathring \R = \R :
\]
kedua operator berulangnya berbeda sejauh yang mungkin bagi \omterm{def:b1:logic:sets}{himpunan}. (iii)
\emph{Gabungan tak hingga \omterm{def:b1:topology:closed}{himpunan tertutup} boleh gagal menjadi \omterm{def:b1:topology:closed}{tertutup}}:
karena $\bigcup_{n\geq1} \intcc{\frac1n}{1} = \intoc{0}{1}$ --- yaitu
cermin contoh penyangkal irisan pada
\cref{prop:b1:topology:openstable}. (iv) \emph{\omterm{def:b1:topology:dense}{Padat} tak
berarti besar}: karena $\Q$ \omterm{def:b1:topology:dense}{padat}, terbilang, dengan \omterm{def:b1:topology:closure}{interior} yang kosong, dan
komplemennya \omterm{def:b1:topology:dense}{padat} pula; jadi \omterm{def:b1:topology:dense}{kepadatan} mengatakan ``sedekat-dekatnya dengan
segalanya'', bukan ``hampir segalanya'' --- adapun \omterm{pb:b1:topology:1}{himpunan Cantor} pada soal akhir
pekan (\cref{pb:b1:topology:1}) menegaskan yang sebaliknya,
yaitu \omterm{def:b1:logic:sets}{himpunan} yang kecil secara topologi namun besar secara takterbilang.
\end{remark}

\begin{example}[Sebuah penutup yang dihitung persis]\label{ex:b1:topology:closurecomputed}
Misalkan $G = \bigl\{\frac1m + \frac1n : m, n \in \N^*\bigr\}$ (dari
\cref{ex:b1:topology:seqtest}). Klaimnya:
\[
\overline G = G \,\cup\, \Bigl\{\frac1m : m \in \N^*\Bigr\}
\,\cup\, \{0\} .
\]
($\supseteq$) Karena $\frac1m = \lim_n \bigl(\frac1m + \frac1n\bigr)$
dan $0 = \lim_n \frac2n$: jadi lekat menurut pencirian lewat
barisan. ($\subseteq$) Misalkan $x = \lim_k \bigl(
\frac{1}{m_k} + \frac{1}{n_k}\bigr)$; lalu urutkan setiap pasangannya sehingga
$m_k \leq n_k$. Jika $(m_k)$ tak terbatas, maka suatu \omterm{def:b1:seq:subsequence}{subbarisan} mempunyai $m_k
\to \infty$, sehingga $n_k \to \infty$ pula dan $x = 0$. Kalau tidak,
$(m_k)$ mengambil nilai yang berhingga banyak, salah satunya, katakanlah $m$,
tak hingga sering; lalu sepanjang \omterm{def:b1:seq:subsequence}{subbarisan} itu $\frac{1}{n_k} \to x -
\frac1m$: jika $(n_k)$ terbatas maka ia mengambil suatu nilai $n$
tak hingga sering dan $x = \frac1m + \frac1n \in G$; kalau tidak, $x =
\frac1m$. Jadi setiap kasusnya mendarat di \omterm{def:b1:logic:sets}{himpunan} yang diumumkan. Inti gagasan
penutupnya: bahwa menghitung \omterm{def:b1:topology:closure}{penutup} adalah analisis kasus bergaya kekompakan
atas indeksnya --- indeks yang terbatas berarti nilai yang berhingga
banyak (menurut sarang merpati), sedangkan indeks yang tak terbatas berarti sebuah limit kabur ---
dan jawabannya memperlihatkan struktur dua lapis yang khas bagi
titik limit: yaitu \omterm{def:b1:logic:sets}{himpunannya}, limit generasi pertamanya, dan
limitnya $0$.
\end{example}

\begin{definition}[Kepadatan, bentuk topologinya]\label{def:b1:topology:dense}
\omterm{def:b1:logic:sets}{Himpunan} $A$ disebut \emph{padat}\index{kepadatan} di $\R$ bila $\overline A = \R$
--- setara dengan itu, bila setiap \omterm{prop:b1:reals:intervals}{selang} \omterm{def:b1:topology:open}{terbuka} yang tak kosong berpotongan dengan $A$;
dan setara pula (menurut \cref{prop:b1:topology:closureprops}~(3)), bila setiap
bilangan real merupakan limit unsur $A$. Contohnya: $\Q$, $\R \setminus
\Q$, bilangan diadik (\cref{exo:b1:reals:8}), dan \omterm{def:b1:structures:subgroup}{subgrup} yang padat
(\cref{exo:b1:reals:9}).
\end{definition}

\begin{example}[Kepadatan bersifat relatif]\label{ex:b1:topology:relativedensity}
``\omterm{def:b1:topology:dense}{Padat}'' sebagaimana didefinisikan di sini berarti \omterm{def:b1:topology:dense}{padat} \emph{di $\R$}; padahal sebuah \omterm{def:b1:logic:sets}{himpunan} boleh
saja \omterm{def:b1:topology:dense}{padat} hanya pada sebagian garisnya. Bilangan diadik pada
$\intcc{0}{1}$, yakni $D \cap \intcc{0}{1}$
(\cref{exo:b1:reals:8}), berpotongan dengan setiap \omterm{prop:b1:reals:intervals}{selang} \omterm{def:b1:topology:open}{terbuka} yang termuat
$\intcc{0}{1}$ tetapi tentu saja melewatkan $\intoo{2}{3}$ seluruhnya: jadi mereka
\omterm{def:b1:topology:dense}{padat} \emph{di} $\intcc{0}{1}$, yang berarti $\overline{D \cap
\intcc{0}{1}} = \intcc{0}{1}$. Adapun frasa umum ``$A$ \omterm{def:b1:topology:dense}{padat}
di $B$'' menyingkat $B \subseteq \overline A$ --- jadi namailah selalu
\omterm{def:b1:logic:sets}{himpunan} lingkungannya, karena titik ujung pada soal akhir pekan bersifat \omterm{def:b1:topology:dense}{padat}
\emph{di \omterm{pb:b1:topology:1}{himpunan Cantor}} sekaligus tak \omterm{def:b1:topology:dense}{padat} di mana pun \emph{di
$\R$}: yaitu \omterm{def:b1:logic:sets}{himpunan} yang sama, dengan dua pemerian yang jujur namun
terdengar berlawanan.
\end{example}

\begin{example}[Menangani kepadatan]\label{ex:b1:topology:densitytricks}
Tiga gerakan cepat yang terus-menerus berulang. \emph{Memperbesar}: jika
$A$ \omterm{def:b1:topology:dense}{padat} dan $A \subseteq B$, maka $B$ \omterm{def:b1:topology:dense}{padat} (karena setiap
\omterm{prop:b1:reals:intervals}{selang} sudah berpotongan dengan $A$). \emph{Mengangkut}: jika $A$
\omterm{def:b1:topology:dense}{padat}, maka $\lambda A + \mu$ pun \omterm{def:b1:topology:dense}{padat} untuk $\lambda \neq 0$ --- karena
\omterm{prop:b1:reals:intervals}{selang} $I$ berpotongan dengan $\lambda A + \mu$ jika dan hanya jika \omterm{prop:b1:reals:intervals}{selang}
$\frac{I - \mu}{\lambda}$ berpotongan dengan $A$; jadi kelipatan ganjil
$10^{-9}$, misalnya, bersifat \omterm{def:b1:topology:dense}{padat}. \emph{Mengiriskan gagal}: karena dua \omterm{def:b1:logic:sets}{himpunan}
yang \omterm{def:b1:topology:dense}{padat} dapat saling melewatkan seluruhnya ($\Q$ dan $\R \setminus \Q$):
jadi \omterm{def:b1:topology:dense}{kepadatan} bertahan pada gabungan dan \omterm{def:b1:logic:map}{pemetaan} afin, tetapi tak pernah pada irisan.
\end{example}

\section{Kekompakan ruas}

\begin{theorem}[Ruas bersifat kompak secara barisan]\label{thm:b1:topology:compact}
Misalkan $a \leq b$. Setiap barisan berisi titik $\intcc{a}{b}$ mempunyai
\omterm{def:b1:seq:subsequence}{subbarisan} yang konvergen ke sebuah titik \emph{pada}
$\intcc{a}{b}$.\index{kekompakan}

Secara lebih umum, \omterm{def:b1:logic:sets}{himpunan} bagian $\R$ yang bersifat demikian (yaitu setiap
barisannya mempunyai \omterm{def:b1:seq:subsequence}{subbarisan} yang konvergen di dalam \omterm{def:b1:logic:sets}{himpunannya}) tepat merupakan
\omterm{def:b1:logic:sets}{himpunan} yang \emph{\omterm{def:b1:topology:closed}{tertutup} dan terbatas}.
\end{theorem}

\begin{proof}
Barisan di $\intcc{a}{b}$ bersifat terbatas, sehingga Bolzano--Weierstrass
(\cref{thm:b1:seq:bw}) mengekstrak \omterm{def:b1:seq:subsequence}{subbarisan} yang konvergen; dan limitnya
tinggal di $\intcc{a}{b}$ karena ruas bersifat \omterm{def:b1:topology:closed}{tertutup}
(\cref{thm:b1:topology:seqclosed}).

Kasus umumnya. \emph{(\omterm{def:b1:topology:closed}{Tertutup} dan terbatas $\Rightarrow$ kompak)}: misalkan
$F$ \omterm{def:b1:topology:closed}{tertutup} dan terbatas, dan $(u_n)$ barisan di $F$.
Keterbatasan $F$ membatasi barisannya, sehingga Bolzano--Weierstrass
mengekstrak $u_{\varphi(n)} \to \ell$; dan $\ell \in F$ karena $F$
\omterm{def:b1:topology:closed}{tertutup} dan \omterm{def:b1:seq:subsequence}{subbarisannya} merupakan barisan yang konvergen berisi titik
$F$ (\cref{thm:b1:topology:seqclosed}): jadi kedua hipotesisnya
terpakai masing-masing satu, keterbatasan untuk keberadaan limitnya,
dan ketertutupan untuk keanggotaannya. \emph{(Kompak $\Rightarrow$ \omterm{def:b1:topology:closed}{tertutup} dan
terbatas)}: jika $F$ tak terbatas, pilihlah $u_n \in F$ dengan $\abs{u_n} \geq
n$; maka setiap \omterm{def:b1:seq:subsequence}{subbarisannya} tak terbatas, sehingga divergen
(\cref{prop:b1:seq:first}): jadi sama sekali tak ada \omterm{def:b1:seq:subsequence}{subbarisan} yang konvergen. Jika $F$
tak \omterm{def:b1:topology:closed}{tertutup}, ambillah $u_n \in F$ dengan $u_n \to \ell \notin F$
(\cref{thm:b1:topology:seqclosed}): maka setiap \omterm{def:b1:seq:subsequence}{subbarisannya} konvergen ke
$\ell \notin F$, sehingga tak ada \omterm{def:b1:seq:subsequence}{subbarisan} yang konvergen \emph{di dalam} $F$.
\end{proof}

\begin{example}[Himpunan kompak yang bersarang]\label{ex:b1:topology:nested}
Sebuah latihan pertama bagi teoremanya. Misalkan $K_0 \supseteq K_1 \supseteq
K_2 \supseteq \dots$ \omterm{def:b1:logic:sets}{himpunan} bagian $\R$ yang kompak (yaitu \omterm{def:b1:topology:closed}{tertutup} dan terbatas)
dan tak kosong. Maka $\bigcap_n K_n \neq \emptyset$. Memang, pilihlah $x_n
\in K_n$ untuk setiap $n$: barisannya hidup di dalam $K_0$ yang kompak,
sehingga suatu \omterm{def:b1:seq:subsequence}{subbarisan} $x_{\varphi(n)}$ konvergen ke suatu $x$
(\cref{thm:b1:topology:compact}). Untuk setiap $m$ yang tetap, suku
$x_{\varphi(n)}$ dengan $\varphi(n) \geq m$ semuanya terletak di \omterm{def:b1:topology:closed}{himpunan
tertutup} $K_m$, sehingga limitnya $x$ terletak di $K_m$
(\cref{thm:b1:topology:seqclosed}); dan karena $m$ tadi sebarang, $x \in
\bigcap_n K_n$. Ada pendamping yang berguna: jika sebuah \omterm{def:b1:topology:open}{himpunan terbuka} $U$ memuat
$\bigcap_n K_n$, maka $U \supseteq K_n$ untuk suatu $n$ --- terapkanlah
argumen yang sama pada titik $x_n \in K_n \setminus U$; karena limitnya
$x$ akan terletak di $\bigcap K_n \subseteq U$, padahal $U$ yang \omterm{def:b1:topology:open}{terbuka} memaksa
$x_{\varphi(n)} \in U$ akhirnya, yang bertentangan. Kedua
\omterm{def:b1:logic:statement}{pernyataannya} gagal tanpa kekompakan: karena $\bigcap_n \intoo{0}{\frac
1n} = \emptyset$ dan $\bigcap_n \intco{n}{+\infty} = \emptyset$.
Inti gagasan penutupnya: bahwa kekompakan mengubah rantai tak hingga atas
penegasan ketakkosongan menjadi satu titik limit --- jadi ia perkakas
yang bertahan pada peralihan ke irisan tak hingga, dan soal
akhir pekan (\cref{pb:b1:topology:1}) akan bersandar padanya dua kali.
\end{example}

\begin{omfigure}
\begin{tikzpicture}[scale=9.5]
  \node[left, font=\small] at (-0.02,0) {$C_0$};
  \fill[omDef] (0,-0.011) rectangle (1,0.011);
  \node[left, font=\small] at (-0.02,-0.09) {$C_1$};
  \fill[omDef] (0,-0.101) rectangle (0.3333,-0.079);
  \fill[omDef] (0.6667,-0.101) rectangle (1,-0.079);
  \node[left, font=\small] at (-0.02,-0.18) {$C_2$};
  \fill[omDef] (0,-0.191) rectangle (0.1111,-0.169);
  \fill[omDef] (0.2222,-0.191) rectangle (0.3333,-0.169);
  \fill[omDef] (0.6667,-0.191) rectangle (0.7778,-0.169);
  \fill[omDef] (0.8889,-0.191) rectangle (1,-0.169);
  \node[left, font=\small] at (-0.02,-0.27) {$C_3$};
  \fill[omDef] (0,-0.281) rectangle (0.0370,-0.259);
  \fill[omDef] (0.0741,-0.281) rectangle (0.1111,-0.259);
  \fill[omDef] (0.2222,-0.281) rectangle (0.2593,-0.259);
  \fill[omDef] (0.2963,-0.281) rectangle (0.3333,-0.259);
  \fill[omDef] (0.6667,-0.281) rectangle (0.7037,-0.259);
  \fill[omDef] (0.7407,-0.281) rectangle (0.7778,-0.259);
  \fill[omDef] (0.8889,-0.281) rectangle (0.9259,-0.259);
  \fill[omDef] (0.9630,-0.281) rectangle (1,-0.259);
  \node[below, font=\small] at (0,-0.30) {$0$};
  \node[below, font=\small] at (0.3333,-0.30) {$\tfrac13$};
  \node[below, font=\small] at (0.6667,-0.30) {$\tfrac23$};
  \node[below, font=\small] at (1,-0.30) {$1$};
  \draw[omThm, thick, -Stealth] (0.5,-0.042) -- (0.5,-0.062);
  \draw[omThm, thick, -Stealth] (0.5,-0.132) -- (0.5,-0.152);
  \draw[omThm, thick, -Stealth] (0.5,-0.222) -- (0.5,-0.242);
\end{tikzpicture}

{\small Empat tahap pertama konstruksi sepertiga tengahnya:
setiap ruas $C_n$ kehilangan sepertiga tengahnya yang \omterm{def:b1:topology:open}{terbuka}, sehingga menyisakan
$2^{n+1}$ ruas $C_{n+1}$, yang panjang totalnya
$(\frac23)^{n+1}$. Adapun \omterm{pb:b1:topology:1}{himpunan Cantor} $C = \bigcap_n C_n$ --- yaitu
pokok bahasan soal akhir pekan \cref{pb:b1:topology:1} --- merupakan
sisa kompak yang tak kosong yang dijamin oleh argumen kompak bersarang
di atas: dengan panjang nol, namun titik yang bertahan
takterbilang banyaknya.}
\end{omfigure}

\begin{remark}
Inilah mesin di balik teorema nilai ekstrem
(\cref{ch:b1:continuity}) dan teorema Heine tentang kekontinuan seragam.
Adapun nama ``kompak'' akan memperoleh definisi umumnya (lewat selimut)
pada tahun kedua; sedangkan pada $\R$, kekompakan barisan sudah semua yang kita perlukan,
dan ``kompak $=$ \omterm{def:b1:topology:closed}{tertutup} $+$ terbatas'' adalah \omterm{def:b1:logic:statement}{pernyataan} yang perlu diingat.
\end{remark}

\begin{example}[Perbatasan di bawah gabungan]\label{ex:b1:topology:boundaryunion}
Selalu berlaku $\partial(A \cup B) \subseteq \partial A \cup \partial
B$: karena sebuah titik $\partial(A \cup B)$ mempunyai setiap \omterm{def:b1:topology:open}{persekitaran}
yang berpotongan dengan $A \cup B$ (sehingga dengan $A$ atau $B$, tak hingga sering salah satu
darinya) dan berpotongan dengan komplemen $A \cup B$, yang terletak di
kedua komplemennya --- lalu pemeriksaan singkat menaruh titiknya di
$\partial A$ atau $\partial B$. Inklusinya dapat tegas secara
mencolok: karena dengan $A = \Q$ dan $B = \R\setminus\Q$,
\[
\partial(A \cup B) = \partial \R = \emptyset ,
\qquad
\partial A \cup \partial B = \R \cup \R = \R :
\]
dua \omterm{def:b1:logic:sets}{himpunan} yang compang-camping dapat merekat menjadi satu yang mulus, dengan \omterm{def:b1:topology:closure}{perbatasannya}
saling memusnahkan. Inti gagasan penutupnya: bahwa \omterm{def:b1:topology:closure}{interior} dan
\omterm{def:b1:topology:closure}{penutup} berperilaku \emph{monoton} di bawah gabungan dan
irisan, tetapi \omterm{def:b1:topology:closure}{perbatasan} tidak --- jadi perlakukanlah $\partial$ sebagai
besaran turunan ($\overline A \setminus \mathring A$), dan jangan pernah sebagai
operator dengan aljabarnya sendiri.
\end{example}

\begin{remark}[Cakrawala di dalam jilid ini]
Kosakata yang dibangun di sini dipakai dua kali lagi dalam buku ini.
Pada \cref{ch:b1:continuity}, setiap teoremanya merupakan \omterm{def:b1:logic:statement}{pernyataan} topologi
yang menyamar: teorema nilai antara mengatakan bahwa
\omterm{def:b1:logic:map}{pemetaan} yang kontinu mengawetkan sifat \omterm{prop:b1:reals:intervals}{selang}, sedangkan teorema nilai
ekstrem mengatakan bahwa ia mengawetkan kekompakan --- dan buktinya
memanggil \cref{thm:b1:topology:seqclosed,thm:b1:topology:compact}
dengan namanya. Pada \cref{ch:b1:multivar}, definisi yang sama
dibaca ulang di $\R^2$ dengan cakram menggantikan \omterm{prop:b1:reals:intervals}{selang}: sehingga \omterm{def:b1:topology:open}{himpunan terbuka},
\omterm{def:b1:topology:closure}{penutup} dan kekompakan terbawa kata demi kata, dan
teorema nilai ekstrem dua variabel kembali menumpang pada
Bolzano--Weierstrass (dengan mengekstrak pada setiap koordinatnya). Adapun satu
gagasan yang \emph{tak} tersamaratakan tanpa nyeri adalah
\omterm{prop:b1:reals:intervals}{selangnya} sendiri --- di bidang, keterhubungan menggantikan
kecembungan, yaitu kisah yang dimulai dengan
``hanya $\emptyset$ dan $\R$ yang \omterm{def:b1:topology:open}{terbuka} sekaligus \omterm{def:b1:topology:closed}{tertutup}'' pada
\cref{exo:b1:topology:9}.
\end{remark}

\section{Latihan}

\begin{exercise}[$\star$]\label{exo:b1:topology:1}
Untuk setiap \omterm{def:b1:logic:sets}{himpunan}, katakanlah apakah ia \omterm{def:b1:topology:open}{terbuka}, \omterm{def:b1:topology:closed}{tertutup}, keduanya, atau bukan keduanya (beserta
pembenarannya): $\intoo{0}{1} \cup \intoo{2}{3}$; $\;\intco{0}{1}$;
$\;\{0\} \cup \intcc{1}{2}$; $\;\R \setminus \Z$; $\;\Q \cap
\intoo{0}{1}$.
\end{exercise}

\begin{exercise}[$\star$]\label{exo:b1:topology:2}
Tentukan $\mathring A$, $\overline A$ dan $\partial A$ untuk:
$A = \intoc{0}{1} \cup \{2\}$; $\;A = \R \setminus \Q$;
$\;A = \bigl\{\frac{(-1)^n n}{n+1} : n \in \N\bigr\}$.
\end{exercise}

\begin{exercise}[$\star$]\label{exo:b1:topology:3}
Buktikan bahwa \omterm{def:b1:counting:card}{himpunan hingga} bersifat \omterm{def:b1:topology:closed}{tertutup}, mula-mula lewat komplemennya, lalu lewat
pencirian barisannya.
\end{exercise}

\begin{exercise}[$\star$]\label{exo:b1:topology:4}
Buktikan bahwa untuk sebarang $A, B \subseteq \R$: $\overline{A \cup B} =
\overline A \cup \overline B$. Tunjukkan lewat contoh bahwa
$\overline{A \cap B}$ dapat berbeda dari $\overline A \cap \overline
B$.
\end{exercise}

\begin{exercise}[$\star\star$]\label{exo:b1:topology:5}
Misalkan $u_n \to \ell$ di $\R$. Buktikan bahwa \omterm{def:b1:logic:sets}{himpunan} $\{u_n : n \in \N\}
\cup \{\ell\}$ bersifat \omterm{def:b1:topology:closed}{tertutup} (sehingga kompak bila kita menambahkan bahwa ia
terbatas --- dan memang demikian).
\end{exercise}

\begin{exercise}[$\star\star$]\label{exo:b1:topology:6}
Misalkan $U$ \omterm{def:b1:topology:open}{terbuka} dan $A$ sebarang. Buktikan bahwa $U + A = \{u + a\}$ bersifat
\omterm{def:b1:topology:open}{terbuka}. Simpulkan bahwa jumlah sebuah \omterm{def:b1:topology:open}{himpunan terbuka} dengan sebarang \omterm{def:b1:logic:sets}{himpunan} bersifat \omterm{def:b1:topology:open}{terbuka}, lalu
bandingkan: tunjukkanlah dua \omterm{def:b1:logic:sets}{himpunan} \emph{\omterm{def:b1:topology:closed}{tertutup}} yang jumlahnya tak \omterm{def:b1:topology:closed}{tertutup}.
\emph{(Cobalah $\Z$ dan $\sqrt 2\,\Z$, dengan \cref{exo:b1:reals:9}.)}
\end{exercise}

\begin{exercise}[$\star\star$]\label{exo:b1:topology:7}
Sebuah titik $x \in A$ disebut \emph{terasing} di $A$ bila suatu \omterm{def:b1:topology:open}{persekitaran}
$x$ berpotongan dengan $A$ hanya pada $x$. Buktikan bahwa setiap titik $\Z$ terasing
di $\Z$, bahwa $A = \{\frac1n\}$ mempunyai semua titiknya terasing namun
$\overline A \neq A$, dan bahwa \omterm{def:b1:logic:sets}{himpunan} yang semua titiknya terasing
mempunyai \omterm{def:b1:topology:closure}{interior} yang kosong.
\end{exercise}

\begin{exercise}[$\star\star$]\label{exo:b1:topology:8}
Buktikan bahwa \omterm{def:b1:topology:closure}{penutup} sebuah \omterm{def:b1:logic:sets}{himpunan} yang terbatas bersifat terbatas, dan bahwa
$\sup \overline A = \sup A$ untuk $A$ yang tak kosong dan terbatas di atas. Simpulkan
bahwa $\sup A \in \overline A$: jadi \omterm{def:b1:reals:bounds}{supremumnya} selalu lekat.
\end{exercise}

\begin{exercise}[$\star\star\star$]\label{exo:b1:topology:9}
Buktikan bahwa satu-satunya \omterm{def:b1:logic:sets}{himpunan} bagian $\R$ yang sekaligus \omterm{def:b1:topology:open}{terbuka} dan \omterm{def:b1:topology:closed}{tertutup} adalah
$\emptyset$ dan $\R$. \emph{Petunjuk: andaikan $A$ \omterm{def:b1:topology:open}{terbuka}, \omterm{def:b1:topology:closed}{tertutup}, dengan
$A \neq \emptyset$ dan $\R \setminus A \neq \emptyset$; pilihlah $a \in
A$, $b \notin A$, katakanlah $a < b$, lalu tinjaulah $s = \sup\,(A \cap
\intcc{a}{b})$; putuskanlah apakah $s$ dapat termasuk $A$ atau
komplemennya.}
\end{exercise}

\begin{exercise}[$\star\star\star$]\label{exo:b1:topology:10}
(Struktur \omterm{def:b1:topology:open}{himpunan terbuka}) Misalkan $U \subseteq \R$ \omterm{def:b1:topology:open}{terbuka} dan tak kosong. Untuk
$x \in U$, misalkan $I_x$ gabungan semua \omterm{prop:b1:reals:intervals}{selang} \omterm{def:b1:topology:open}{terbuka} yang memuat
$x$ dan termuat $U$. Buktikan bahwa $I_x$ merupakan \omterm{prop:b1:reals:intervals}{selang} \omterm{def:b1:topology:open}{terbuka}, bahwa
dua \omterm{def:b1:logic:sets}{himpunan} $I_x$, $I_y$ sama atau saling lepas, dan bahwa $U$ merupakan gabungan
atas \omterm{prop:b1:reals:intervals}{selang} \omterm{def:b1:topology:open}{terbuka} yang saling lepas dan \emph{terbilang banyaknya} \emph{(pilihlah
satu bilangan rasional pada masing-masingnya)}.
\end{exercise}

\begin{exercise}[$\star\star$]\label{exo:b1:topology:11}
Sebuah titik $x \in \R$ disebut \emph{titik akumulasi} $A$ bila
setiap \omterm{def:b1:topology:open}{persekitaran} $x$ berpotongan dengan $A \setminus \{x\}$; dan \omterm{def:b1:logic:sets}{himpunannya} adalah
\emph{\omterm{def:b1:logic:sets}{himpunan} turunan} $A'$. Buktikan bahwa $\overline A = A \cup A'$,
dan bahwa $A$ \omterm{def:b1:topology:closed}{tertutup} jika dan hanya jika $A' \subseteq A$. Tentukan
$A'$ untuk $A = \{\frac 1n : n \in \N^*\}$, untuk $A = \Z$, dan untuk
$A = \Q$.
\end{exercise}

\begin{exercise}[$\star\star\star$]\label{exo:b1:topology:12}
Untuk $A \subseteq \R$ yang tak kosong, definisikan $d_A(x) = \inf\{\abs{x - a}
: a \in A\}$. Buktikan:
\begin{enumerate}
  \item $\abs{d_A(x) - d_A(y)} \leq \abs{x - y}$ untuk setiap $x, y$
        (jadi $d_A$ bersifat Lipschitz-$1$);
  \item $d_A(x) = 0$ jika dan hanya jika $x \in \overline A$;
        khususnya, jika $F$ \omterm{def:b1:topology:closed}{tertutup} dan $x \notin F$, maka
        $d_F(x) > 0$;
  \item untuk setiap $\varepsilon > 0$, \omterm{def:b1:logic:sets}{himpunan} $V_\varepsilon =
        \{x : d_F(x) < \varepsilon\}$ bersifat \omterm{def:b1:topology:open}{terbuka}, memuat $F$, dan
        $\bigcap_{\varepsilon > 0} V_\varepsilon = F$ untuk $F$ yang
        \omterm{def:b1:topology:closed}{tertutup}: jadi setiap \omterm{def:b1:topology:closed}{himpunan tertutup} merupakan irisan terbilang atas
        \omterm{def:b1:topology:open}{himpunan terbuka}.
\end{enumerate}
\end{exercise}

\section{Soal: Himpunan Cantor, kecil dan raksasa sekaligus}

\begin{problem}[{Soal akhir pekan --- \omterm{pb:b1:topology:1}{himpunan Cantor} sepertiga
tengah: panjang nol, takterbilang, sempurna, dan $C + C =
\intcc{0}{2}$}]
\label{pb:b1:topology:1}
Buanglah dari $\intcc{0}{1}$ sepertiga tengahnya yang \omterm{def:b1:topology:open}{terbuka}, lalu sepertiga
tengah setiap ruas yang tersisa, dan ulangilah selamanya: yang
bertahan adalah \emph{\omterm{pb:b1:topology:1}{himpunan Cantor}} $C$, yaitu pabrik
contoh penyangkal yang mendasar dalam analisis. Soal ini membangunnya,
membacanya lewat mesin basis $3$ pada \cref{pb:b1:reals:1},
lalu menegakkan potret paradoksalnya: panjang totalnya nol, namun
takterbilang; \omterm{def:b1:topology:closure}{interiornya} kosong, namun tanpa titik yang terasing; terputus
seluruhnya, namun $C + C$ mengisi seluruh ruas
$\intcc{0}{2}$.\index{himpunan Cantor} Secara resmi: $C_0 = \intcc{0}{1}$,
dan $C_{n+1}$ diperoleh dari $C_n$ dengan menghapus sepertiga tengah
yang \omterm{def:b1:topology:open}{terbuka} dari setiap ruas $C_n$; akhirnya $C = \bigcap_{n \geq 0}
C_n$. Di sepanjang soal ini, \emph{sandi terner} bagi $x \in \intcc{0}{1}$
adalah sebarang untaian angka $(d_k)_{k\geq1}$ dengan $d_k \in \{0, 1, 2\}$
yang nilainya $\sup_n \sum_{k=1}^n d_k 3^{-k}$ sama dengan $x$ ---
adapun sandi yang tak sejati (yang akhirnya $2$) diperbolehkan; dan menurut
\cref{pb:b1:reals:1} (pertanyaan 9--11), setiap $x \in
\intcc{0}{1}$ mempunyai satu atau dua sandi, dan dua tepat ketika $x = m/3^N
\in \intoo{0}{1}$.

\medskip
\noindent\textbf{Bagian I --- Konstruksinya.}
\begin{enumerate}
  \item Perikan $C_1$ dan $C_2$ secara eksplisit sebagai gabungan
        ruas, lalu buktikan lewat induksi: bahwa $C_n$ merupakan gabungan
        lepas atas $2^n$ ruas \omterm{def:b1:topology:closed}{tertutup}, yang masing-masingnya berpanjang $3^{-n}$.
  \item Tunjukkan bahwa $C$ \omterm{def:b1:topology:closed}{tertutup}, terbatas --- sehingga kompak
        (\cref{thm:b1:topology:compact}) --- tak kosong, dan bahwa
        setiap titik ujung setiap ruas setiap $C_n$ termasuk
        $C$.
  \item Panjang total $C_n$ adalah $\bigl(\frac23\bigr)^n$.
        Simpulkan bahwa untuk setiap $\varepsilon > 0$, \omterm{def:b1:logic:sets}{himpunan} $C$ dapat
        diselimuti oleh ruas yang berhingga banyak dengan panjang total
        $\leq \varepsilon$: jadi \omterm{pb:b1:topology:1}{himpunan Cantor} berpanjang \emph{nol}.
  \item Tunjukkan bahwa \omterm{prop:b1:reals:intervals}{selang} yang termuat $C$ berpanjang $\leq
        3^{-n}$ untuk setiap $n$, sehingga ia singleton atau kosong:
        jadi $\mathring C = \emptyset$. Karena \omterm{def:b1:topology:closed}{tertutup} dengan \omterm{def:b1:topology:closure}{interior}
        yang kosong, $C$ bersifat \emph{tak \omterm{def:b1:topology:dense}{padat} di mana pun}.
\end{enumerate}

\medskip
\noindent\textbf{Bagian II --- Sandi ternernya.}
\begin{enumerate}[resume]
  \item Buktikan rekursi keserupaan dirinya
        \[
        C_{n+1} = \tfrac13 C_n \,\cup\,
        \bigl(\tfrac23 + \tfrac13 C_n\bigr),
        \qquad\text{sehingga}\qquad
        C = \tfrac13 C \,\cup\, \bigl(\tfrac23 + \tfrac13
        C\bigr),
        \]
        dengan kedua kepingannya saling lepas: jadi $C$ adalah dua salinan
        dirinya sendiri pada skala $\frac13$.
  \item Buktikan lewat induksi pada $n$: bahwa $x \in C_n$ jika dan hanya jika
        $x$ mempunyai sandi terner yang $n$ angka pertamanya terletak di
        $\{0, 2\}$. Lalu simpulkan, dengan memakai kenyataan bahwa $x$ mempunyai paling banyak
        dua sandi: bahwa $x \in C$ jika dan hanya jika $x$ mempunyai sandi
        \emph{tanpa angka yang sama dengan $1$} (yaitu sandi \emph{bebas-$1$}).
  \item Sandi dalam kerja: berikanlah sandi bebas-$1$ bagi $0$, $1$,
        $\frac13$, $\frac23$; tunjukkan $\frac14 = (0.\overline{02})_3$
        dan $\frac34 = (0.\overline{20})_3$, sehingga keduanya termasuk
        $C$; lalu periksalah bahwa $\frac14$ \emph{bukan} titik ujung
        $C_n$ mana pun (karena titik ujung berbentuk $m/3^n$).
  \item Tunjukkan bahwa setiap $x \in C$ mempunyai \emph{tepat satu} sandi
        bebas-$1$ \emph{(ketika $x$ mempunyai dua sandi, buktikan bahwa tepat
        satu dari pasangannya memuat angka $1$)}. Simpulkan: bahwa
        \omterm{def:b1:logic:map}{pemetaan} nilainya merupakan bijeksi dari untaian $\{0,2\}$ pada
        $C$.
  \item (Diagonal) Misalkan $k \mapsto x_k$ sebarang \omterm{def:b1:logic:map}{pemetaan} $\N^* \to C$.
        Bangunlah untaian $\{0, 2\}$ yang berbeda pada indeks $k$ dari
        sandi $x_k$, lalu simpulkan bahwa $C$ takterbilang
        --- sedangkan sebaliknya, pertanyaan 3 mengatakan bahwa ia
        terabaikan secara metrik.
\end{enumerate}

\medskip
\noindent\textbf{Bagian III --- Potret topologinya.}
\begin{enumerate}[resume]
  \item Rakitlah catatan sejauh ini: bahwa $C$ kompak, takterbilang,
        berpanjang nol, dan tak \omterm{def:b1:topology:dense}{padat} di mana pun. Pemuatan tunggal
        $C \subseteq C_n$ yang mana yang membawa masing-masing sifatnya?
  \item ($C$ bersifat sempurna) Misalkan $x \in C$ dengan sandi bebas-$1$
        $(d_k)$. Membalikkan angkanya $d_n$ ($0 \leftrightarrow
        2$) menghasilkan $x_n \in C$ dengan $\abs{x_n - x} = 2 \cdot
        3^{-n}$. Simpulkan bahwa $C$ tak mempunyai titik terasing: bahwa setiap
        titik $C$ merupakan limit titik \emph{lain} pada $C$.
  \item Tunjukkan bahwa titik ujung pada pertanyaan 2 membentuk \omterm{def:b1:logic:sets}{himpunan} bagian
        $C$ yang terbilang dan \omterm{def:b1:topology:dense}{padat} \emph{(pengallah sandinya setelah $n$
        angka lalu lanjutkan dengan $0$; adapun keterbilangannya seperti pada
        \cref{exo:b1:topology:10})}. Simpulkan: bahwa titik yang khas
        pada $C$ --- seperti $\frac14$ --- \emph{bukan} titik ujung:
        jadi titik ujung merupakan kerangka terbilang di dalam tubuh yang
        takterbilang.
  \item (Terputus seluruhnya) Misalkan $x < y$ di $C$. Pilihlah $n$
        dengan $3^{-n} < y - x$ lalu hasilkan sebuah titik $z \in
        \intoo{x}{y}$ dengan $z \notin C$. Simpulkan bahwa satu-satunya
        \omterm{prop:b1:reals:intervals}{selang} tak kosong yang termuat $C$ adalah singleton.
\end{enumerate}

\medskip
\noindent\textbf{Bagian IV --- Aritmetika $C$.}
\begin{enumerate}[resume]
  \item Tunjukkan $1 - C = C$ \emph{(apa yang dilakukan $x \mapsto 1 - x$
        terhadap sandi bebas-$1$? ingatlah $1 = (0.\overline{2})_3$)}.
  \item (Penjumlahan sandi) Tunjukkan bahwa jika $x$, $x'$ bersandi
        $(a_k)$, $(b_k)$, maka $x + x' = \lim_n\,(t_n + t'_n)$
        dengan $t_n, t'_n$ jumlah parsialnya. Simpulkan: bahwa setiap $y
        \in \intcc{0}{1}$ merupakan \emph{titik tengah} dua titik
        $C$ --- karena diberikan sandi $(e_k)$ bagi $y$, pilihlah angka
        $a_k, b_k \in \{0, 2\}$ dengan $\frac{a_k + b_k}{2} =
        e_k$.
  \item Simpulkan $C + C = \intcc{0}{2}$ dan, dengan pertanyaan 14,
        $C - C = \intcc{-1}{1}$. Contoh konkretnya: tulislah $1$ sebagai
        jumlah dua bukan titik ujung yang ditemukan pada pertanyaan 7.
  \item Renungkanlah: sebuah \omterm{def:b1:logic:sets}{himpunan} berpanjang nol yang \omterm{def:b1:logic:sets}{himpunan} selisihnya mengisi
        $\intcc{-1}{1}$. Mengapa tak ada pertentangan antara
        ``$C$ terabaikan secara metrik'' dan ``$C + C$ berpanjang
        penuh''? (Satu kalimat; pikirkanlah apa yang dikendalikan panjang
        dan apa yang tidak.)
\end{enumerate}

\medskip
\noindent\textbf{Bagian V --- Anggotanya, yang rasional dan yang irasional.}
\begin{enumerate}[resume]
  \item Gabungkan pertanyaan 8 dengan kriteria keperiodikan pada
        \cref{pb:b1:reals:1}: bahwa titik $C$ bersifat rasional jika dan
        hanya jika sandi bebas-$1$-nya periodik pada akhirnya. Jalankanlah
        pembagian panjang basis $3$ untuk memeriksa $\frac1{13} =
        (0.\overline{002})_3 \in C$.
  \item Hasilkan sebuah anggota $C$ yang \emph{irasional} secara eksplisit:
        yaitu nilai sandi dengan $d_k = 2$ pada posisi
        segitiga $k = \frac{j(j+1)}{2}$ dan $d_k = 0$
        di tempat lain. Benarkanlah keirasionalannya lewat argumen jarak yang
        tumbuh pada \cref{pb:b1:reals:1} (pertanyaan 20).
  \item (Pada seluruh ruasnya) Tinjaulah $h$ yang memetakan titik
        $C$ bersandi bebas-$1$ $(d_k)$ ke nilai untaian
        \emph{biner} $\bigl(\frac{d_k}2\bigr)$, yakni
        $h(x) = \sup_n \sum_{k=1}^n \frac{d_k}{2}\,2^{-k}$.
        Tunjukkan bahwa $h$ memetakan $C$ \emph{pada} $\intcc{0}{1}$. Jadi
        $C$ yang terabaikan itu \omterm{def:b1:logic:inj}{surjektif} pada sebuah ruas yang berpanjang
        penuh --- yaitu bukti kedua bahwa $C$ takterbilang.
  \item (Panjang yang serupa diri) Andaikan suatu gagasan panjang $L$
        terdefinisi bagi $C$ dan salinan susutannya, yang menghormati
        penskalaan ($L(\lambda A) = \lambda L(A)$), kekekalan terhadap translasi,
        dan keaditifan atas penguraian lepas
        pada pertanyaan 5. Tunjukkan bahwa lalu $L(C) = \frac23\,L(C)$,
        yang memaksa $L(C) = 0$: jadi keserupaan diri saja sudah
        menghukum $C$ menjadi berpanjang nol.
\end{enumerate}

\medskip
\noindent\textbf{Bagian VI --- Sepupunya yang gemuk, dan moralnya.}
\begin{enumerate}[resume]
  \item (\omterm{pb:b1:topology:1}{Himpunan Cantor} yang gemuk) Ulangilah konstruksinya, tetapi pada tahap
        $n$ ($n = 0, 1, 2, \dots$) buanglah dari masing-masing $2^n$
        ruas yang ada sebuah \omterm{prop:b1:reals:intervals}{selang} \omterm{def:b1:topology:open}{terbuka} di tengahnya yang berpanjang
        $4^{-(n+1)}$ saja. Tunjukkan bahwa panjang ruasnya $l_n$
        menuruti $l_{n+1} = \frac{l_n - 4^{-(n+1)}}{2}$, $l_n =
        \frac{2^n + 1}{2\cdot 4^n} > 0$, bahwa hasilnya
        $K = \bigcap K_n$ kompak dengan \omterm{def:b1:topology:closure}{interior} yang kosong, dan
        bahwa panjang total yang terbuang adalah $\sum_{n\geq0} 2^n
        4^{-(n+1)} = \frac12$. Dengan menerima keaditifan panjang bagi
        gabungan hingga atas \omterm{prop:b1:reals:intervals}{selang} (yang intuitif, dan dibahas pada Tahun ke-3), lalu
        memakai kedua \omterm{def:b1:logic:statement}{pernyataan} \cref{ex:b1:topology:nested},
        tunjukkan bahwa sebarang keluarga hingga \omterm{prop:b1:reals:intervals}{selang} \omterm{def:b1:topology:open}{terbuka} yang menyelimuti
        $K$ berpanjang total $\geq \frac12$: jadi $K$ tak \omterm{def:b1:topology:dense}{padat}
        di mana pun tetapi \emph{tak} terabaikan. Jadi kekecilan mempunyai beberapa
        makna yang tak setara.
  \item (Jarak) Tunjukkan bahwa untuk $F \subseteq \R$ yang \omterm{def:b1:topology:closed}{tertutup} dan tak kosong
        serta $x \in \R$, \omterm{def:b1:reals:bounds}{infimum} $d(x, F)$ bersifat
        \emph{tercapai} \emph{(lewat barisan peminimum ditambah
        Bolzano--Weierstrass)}. Lalu hitunglah
        \[
        \max_{y \in \intcc{0}{1}} d(y, C) = \frac16 ,
        \]
        yang tercapai persis di pusatnya $y = \frac12$ \emph{(karena
        titik pada celah yang tercipta pada tahap $n$ berada dalam jarak
        $\frac{3^{-n}}{2}$ dari titik ujung celahnya, yang terletak di
        $C$)}.
  \item (Setiap titik menjadi limit \omterm{def:b1:seq:subsequence}{subbarisan}) Dengan memakai pertanyaan 12
        dan 2, hasilkan satu barisan tunggal di $C$ yang \omterm{def:b1:logic:sets}{himpunan}
        limit \omterm{def:b1:seq:subsequence}{subbarisannya} adalah \emph{seluruh} $C$. (Bandingkanlah:
        bagi barisan yang konvergen \omterm{def:b1:logic:sets}{himpunan} itu satu titik ---
        sedangkan $C$ mewujudkan ujung yang berlawanan di antara yang kompak.)
  \item Sintesis, satu kalimat untuk masing-masing: (i) teorema bab ini
        yang mana yang benar-benar dipakai konstruksinya (kestabilan
        \omterm{def:b1:topology:closed}{himpunan tertutup}, kekompakan, dan pencirian
        lewat barisan)? (ii) daftarkanlah keempat pasangan paradoksal
        pada potretnya (panjang nol/takterbilang,
        \omterm{def:b1:topology:closed}{tertutup}/\omterm{def:b1:topology:closure}{interior} kosong, sempurna/terputus seluruhnya,
        terabaikan/$C+C$ penuh); (iii) di manakah $C$ muncul kembali
        nanti (yaitu tangga setan yang dibangun di atas $h$ pada teori
        kekontinuan, dan teori ukuran pada jilid Tahun ke-3,
        yang di situ $C$ memisahkan ``terbilang'' dari
        ``terabaikan'')?

\end{enumerate}
\end{problem}
